\section{Introduction}
\label{sec:intro}

Continual learning aims to learn a sequence of tasks without catastrophically
forgetting previously acquired knowledge~\cite{french1999catastrophic}. The dominant
strategies --- rehearsal~\cite{rebuffi2017icarl}, regularization~\cite{kirkpatrick2017overcoming,zenke2017continual},
and dynamic architectures~\cite{rusu2016progressive} --- were largely developed in the
regime of training a model from scratch. Recently, however, the community has shifted
toward \emph{pre-trained models} (e.g., Vision Transformers) combined with
\emph{parameter-efficient fine-tuning} (PEFT), in which the backbone is frozen and only a
tiny set of parameters is learned per task~\cite{wang2022l2p,smith2023coda,wang2023olora}.

Low-rank adaptation (LoRA)~\cite{hu2022lora} is a particularly attractive PEFT primitive
for continual learning: a shared LoRA adapter keeps the total parameter count fixed
regardless of the number of tasks, which is ideal for resource-constrained deployment.
However, sequentially fine-tuning a shared adapter still suffers from catastrophic
forgetting --- although the number of trainable parameters is small, the low-rank
subspace is repeatedly overwritten, and the representational structure learned for
earlier tasks is destroyed.

A line of recent work addresses this by allocating \emph{orthogonal} low-rank subspaces
to different tasks~\cite{wang2023olora,liang2024inflora}. These methods prevent
interference by making the new task's adapter (nearly) orthogonal to those of previous
tasks. While effective, they exhibit two limitations. First, they protect all directions
of the past subspace \emph{equally}, ignoring the fact that different directions carry
vastly different amounts of information about a task. Second, most of them allocate a
\emph{new} adapter per task, so the parameter count grows linearly with the number of
tasks.

In this paper we propose \textbf{FOLoRA} (Fisher-Orthogonalized Low-Rank Adaptation), a
replay-free, fixed-budget method that mitigates forgetting by protecting past tasks'
\emph{important} directions in proportion to their importance. Our contributions are:

\begin{enumerate}
\item We frame forgetting in low-rank continual learning through a second-order lens and
derive an upper bound that decomposes forgetting into Fisher-weighted subspace overlaps
(Sec.~\ref{sec:theory}).
\item We propose FOLoRA, which estimates a per-layer second-order importance kernel from
the empirical Fisher information restricted to the adapter parameters, forms an
importance-weighted protected subspace from its top eigendirections, and regularizes the
adapter to stay orthogonal to these directions (Sec.~\ref{sec:method}).
\item We show that FOLoRA's regularizer is exactly the Lagrangian relaxation of the
derived bound, realizing a principled stability--plasticity trade-off, and that it
subsumes equal-weight orthogonalization (O-LoRA) as a special case.
\item On Split CIFAR-100 and Split ImageNet-R, FOLoRA consistently
outperforms sequential LoRA, EWC-LoRA, O-LoRA, L2P and CODA-Prompt in both average
accuracy and forgetting, with no replay buffer and no growth in parameters
(Sec.~\ref{sec:experiments}).
\end{enumerate}
